\section*{Capítulo \ref{ch:g11:ocean-ridges} --- Las dorsales oceánicas y la litosfera oceánica}

\begin{solution}{exo:g11:ocean-ridges:1}
La \omterm{def:g11:ocean-ridges:peridotite}{peridotita} es una roca oscura, de grano grueso, formada sobre todo por
olivino, con algo de piroxeno; gran parte del manto está hecha de ella. El
\omterm{def:g11:ocean-ridges:gabbro}{gabro} es una roca ígnea oscura, de grano grueso, de feldespato plagioclasa y
piroxeno, la roca más abundante de la \omterm{def:g11:ocean-ridges:oceanic-lithosphere}{corteza oceánica} profunda. El \omterm{def:g11:ocean-ridges:gabbro}{gabro}
pertenece a la corteza, y la \omterm{def:g11:ocean-ridges:peridotite}{peridotita}, al manto.
El \omterm{def:g11:ocean-ridges:gabbro}{gabro} cristalizó lentamente, a gran profundidad bajo el fondo marino, donde
el magma se enfrió despacio: sus cristales tuvieron tiempo de crecer. La lava
basáltica se enfrió deprisa en el fondo marino, y sus cristales siguieron
siendo pequeños.
\end{solution}

\begin{solution}{exo:g11:ocean-ridges:2}
Desde el fondo marino hacia abajo: sedimentos (ninguno en el eje), lavas
basálticas (almohadillas y coladas tabulares), complejo de diques laminares,
\omterm{def:g11:ocean-ridges:gabbro}{gabro} y luego, bajo el \omterm{def:g11:seismic-structure:moho}{Moho}, la \omterm{def:g11:ocean-ridges:peridotite}{peridotita} del manto. El sondeo atravesó
$1507/7100 \approx 0.21$, alrededor de una quinta parte de la corteza: ninguna
perforación ha alcanzado nunca los \omterm{def:g11:ocean-ridges:gabbro}{gabros} profundos ni el \omterm{def:g11:seismic-structure:moho}{Moho}.
\end{solution}

\begin{solution}{exo:g11:ocean-ridges:3}
$d(0) = \qty{2500}{m}$; $d(9) = 2500 + 350 \times 3 = \qty{3550}{m}$;
$d(25) = 2500 + 350 \times 5 = \qty{4250}{m}$.
\end{solution}

\begin{solution}{exo:g11:ocean-ridges:4}
Una \omterm{def:g11:ocean-ridges:normal-fault}{falla normal} es una falla inclinada a lo largo de la cual el bloque
situado sobre el plano de falla ha descendido respecto del bloque de debajo.
En una dorsal lenta, las placas tiran de la \omterm{def:g11:lithosphere-asthenosphere:lithosphere}{litosfera}, fría y frágil, en
sentidos opuestos: se estira y se rompe. Los bloques situados sobre las fallas
se deslizan hacia abajo, hacia el eje, y el valle axial es el bloque que se ha
hundido entre ellas.
\end{solution}

\begin{solution}{exo:g11:ocean-ridges:5}
No hierve a causa de la presión: sobre ella pesan más de dos kilómetros de
agua, y a esa presión el agua sigue siendo líquida muy por encima de
\qty{100}{\celsius}. El ``humo'' negro es una nube de granos diminutos de
sulfuros de hierro y otros minerales, que precipitan cuando el agua caliente
de la fuente, rica en metales disueltos, se encuentra con el agua de mar fría.
\end{solution}

\begin{solution}{exo:g11:ocean-ridges:6}
(1) A unos \qty{54}{km} de profundidad, a unos \qty{1324}{\celsius}. (2) A
unos \qty{85}{km} de profundidad, a unos $1400 + 0.45 \times 85 \approx
\qty{1438}{\celsius}$. (3) El manto más caliente empieza a fundirse a más
profundidad, así que se funde una columna más larga de roca ascendente antes
de llegar a la corteza: forma más basalto, y una corteza más gruesa.
\end{solution}

\begin{solution}{exo:g11:ocean-ridges:7}
(1) $1489 - 1345 = \qty{144}{\celsius}$ por debajo de su temperatura de
fusión. (2) Sube $100 - 54 = \qty{46}{km}$; la diferencia se cierra
$144/46 \approx \qty{3.1}{\celsius}$ por kilómetro de media. (3) No: durante el
ascenso la roca incluso se enfría un poco, unos $0.45 \times 46 \approx
\qty{21}{\celsius}$. Lo que baja, mucho más deprisa, es su temperatura de
fusión, a medida que disminuye la presión.
\end{solution}

\begin{solution}{exo:g11:ocean-ridges:8}
Las semivelocidades son \qty{7.3}{cm} y \qty{1.65}{cm} al año, es decir,
\qty{73}{km} y \qty{16.5}{km} por millón de años. Para formar \qty{1000}{km}
de fondo a un lado: $1000/73 \approx 13.7$ millones de años en la dorsal del
Pacífico oriental, y $1000/16.5 \approx 61$ millones de años en la dorsal
mesoatlántica.
\end{solution}

\begin{solution}{exo:g11:ocean-ridges:9}
(1) \qty{5}{cm} son \qty{5e-5}{km}: $65000 \times 5 \times 10^{-5} =
\qty{3.25}{km^2}$ de fondo nuevo al año. (2) $3.25 \times 7.1 \approx
\qty{23}{km^3}$ de corteza al año. (3) $3.6 \times 10^8 / 3.25 \approx 1.1
\times 10^8$ años, unos 110 millones de años: a este ritmo, las dorsales
podrían formar todo el fondo oceánico en unos cien millones de años, poco
tiempo comparado con la edad de la Tierra. Por eso ningún fondo oceánico es
muy antiguo.
\end{solution}

\begin{solution}{exo:g11:ocean-ridges:10}
(1) El plano de falla forma \ang{60} con la horizontal, y el hundimiento es el
cateto vertical de un triángulo rectángulo cuya hipotenusa es el
desplazamiento: $1.5/\sin\ang{60} \approx \qty{1.73}{km}$. (2) Cada falla
separa horizontalmente los lados $1.5/\tan\ang{60} \approx \qty{0.87}{km}$;
las dos fallas juntas, unos \qty{1.73}{km}. (3) No. Los \qty{25}{km} son la
distancia entre las fallas que limitan el bloque hundido, no la cantidad de
estiramiento, que es solo de unos \qty{1.7}{km}; además, la dorsal sigue
abriéndose por nuevos diques y lavas en el eje, no solo por fallamiento.
\end{solution}

\begin{solution}{exo:g11:ocean-ridges:11}
$d(25) = 6400 - 3200\,e^{-25/62.8} \approx 6400 - 3200 \times 0.672 \approx
\qty{4251}{m}$; $d(100) = 6400 - 3200\,e^{-100/62.8} \approx 6400 - 3200
\times 0.203 \approx \qty{5749}{m}$. Entre 25 y 100 millones de años, el fondo
se hunde unos \qty{1500}{m}. La ley de la raíz cuadrada da
$2500 + 350 \times 5 = \qty{4250}{m}$ a los 25 millones de años: las dos
leyes concuerdan donde ambas se aplican.
\end{solution}

\begin{solution}{exo:g11:ocean-ridges:12}
(1) $\alpha\,\Delta T\,L = 3.2 \times 10^{-5} \times 700 \times 100 =
\qty{2.24}{km}$. (2) $s = 3380 \times 2.24/(3380 - 1025) = 7571/2355 \approx
\qty{3.2}{km}$. (3) La ley exponencial permite en total $6400 - 2500 =
\qty{3900}{m}$ de hundimiento: el enfriamiento y la contracción por sí solos
explican unas cuatro quintas partes, con valores sencillos para el descenso de
temperatura y el espesor.
\end{solution}

\begin{solution}{exo:g11:ocean-ridges:13}
(1) $54 - 7 = \qty{47}{km}$ frente a $85 - 7 = \qty{78}{km}$: la columna más
caliente es $78/47 \approx 1.66$ veces más larga. (2) $7.1 \times 1.66 \approx
\qty{11.8}{km}$. (3) Los \qty{10.3}{km} medidos superan lo normal, como se
predice, pero quedan por debajo de la estimación: los kilómetros más profundos
de la columna, justo por encima del inicio de la fusión, dan menos fundido que
los más someros, así que la hipótesis del mismo fundido por kilómetro
sobrestima.
\end{solution}

\begin{solution}{exo:g11:ocean-ridges:14}
(1) La ley da \qty{3178}{m}, \qty{4160}{m}, \qty{4782}{m} y
\qty{5267}{m}; las medianas medidas son menos hondas en \qty{64}{m},
\qty{142}{m}, \qty{317}{m} y \qty{619}{m}. (2) La diferencia crece con la
edad. (3) La ley da la profundidad del propio fondo, corregida de sus
sedimentos; las celdas medidas conservan sus sedimentos, más gruesos sobre el
fondo más antiguo, e incluyen montes submarinos y mesetas volcánicas, menos
hondos que el fondo que los rodea.
\end{solution}

\begin{solution}{exo:g11:ocean-ridges:15}
(1) $p = 1025 \times 9.8 \times 2500 \approx \qty{25.1}{MPa}$, más
\qty{0.1}{MPa}: unos \qty{25.2}{MPa}. (2) Está por encima de la presión
crítica, \qty{22.064}{MPa}: el agua de la fuente no puede hervir, por caliente
que esté. (3) La presión alcanza \qty{22.064}{MPa} a la profundidad
$(22.064 - 0.1) \times 10^6/(1025 \times 9.8) \approx \qty{2190}{m}$: ninguna
fuente situada a más de unos \qty{2200}{m} de profundidad puede hervir.
\end{solution}

\begin{solution}{pb:g11:ocean-ridges:1}
\textbf{1.} $14.6/2 = \qty{7.3}{cm}$ al año, es decir, \qty{73}{km} por millón
de años.

\textbf{2.} A la profundidad del eje de la dorsal, $d(0) = \qty{2500}{m}$.

\textbf{3.} Desde el fondo marino hacia abajo: lavas basálticas (almohadillas
y coladas tabulares), complejo de diques laminares y \omterm{def:g11:ocean-ridges:gabbro}{gabro}, unos \qty{7.1}{km}
de corteza en total; bajo el \omterm{def:g11:seismic-structure:moho}{Moho}, la \omterm{def:g11:ocean-ridges:peridotite}{peridotita} del manto, hasta la base de
la \omterm{def:g11:lithosphere-asthenosphere:lithosphere}{litosfera}, que en el eje es todavía muy delgada. Más tarde se depositarán
sedimentos sobre él.

\textbf{4.} Cuando las placas se separan, la \omterm{def:g11:lithosphere-asthenosphere:lithosphere}{astenosfera} asciende bajo la
dorsal. Su temperatura solo baja unos \qty{0.45}{\celsius} por kilómetro,
mientras que su temperatura de fusión baja mucho más deprisa a medida que
disminuye la presión. A unos \qty{54}{km} de profundidad las dos se igualan:
la roca empieza a fundirse, y sigue fundiéndose mientras asciende, sin que se
le aporte calor (\omterm{def:g11:ocean-ridges:partial-melting}{fusión por descompresión}).

\textbf{5.} Cada millón de años, la dorsal forma $2 \times 73 = \qty{146}{km}$
de fondo a sus dos lados; por cada kilómetro de dorsal, el volumen de corteza
es $146 \times 7.1 \times 1 \approx \qty{1040}{km^3}$.

\textbf{6.} $73 \times t$: \qty{73}{km}, \qty{292}{km}, \qty{1168}{km} y
\qty{4672}{km}.

\textbf{7.} $2000/73 \approx 27$ millones de años.

\textbf{8.} Las placas cambian de velocidad y de dirección a lo largo de
decenas de millones de años, como muestra el codo de la cadena
Hawái--Emperador (el \cref{ch:g11:plate-motions}); la velocidad de hoy no
tiene por qué mantenerse durante 64 millones de años.

\textbf{9.} Lejos del eje, la corteza ya no está sobre una lente de magma: se
enfría, y ya no hay debajo calor suficiente para mantener la circulación y
calentar el agua.

\textbf{10.} Al oeste de la dorsal del Pacífico oriental, el trozo viaja sobre
la placa del Pacífico. Hacia el oeste, la placa del Pacífico termina a lo
largo de fosas profundas, donde se hunde bajo sus vecinas: zonas de subducción (el \cref{ch:g11:subduction-zones}).

\textbf{11.} $d(1) = \qty{2850}{m}$, $d(4) = 2500 + 700 = \qty{3200}{m}$,
$d(16) = 2500 + 1400 = \qty{3900}{m}$, $d(64) = 2500 + 2800 = \qty{5300}{m}$.

\textbf{12.} $3900 - 2500 = \qty{1400}{m}$ durante los primeros 16 millones de
años, y $5300 - 3900 = \qty{1400}{m}$ entre 16 y 64 millones de años: el mismo
hundimiento, en un tiempo tres veces más largo. El fondo se hunde cada vez más
despacio a medida que envejece.

\textbf{13.} $350\sqrt{t} = 1500$ da $\sqrt{t} \approx 4.29$, luego
$t \approx 18.4$ millones de años.

\textbf{14.} $6400 - 3200\,e^{-64/62.8} \approx 6400 - 3200 \times 0.361
\approx \qty{5245}{m}$, unos \qty{55}{m} menos que la ley de la raíz cuadrada:
las dos leyes concuerdan con un margen de alrededor del \qty{1}{\percent}.

\textbf{15.} La mediana medida queda unos \qty{650}{m} menos honda que la ley.
La ley describe el fondo corregido del peso de sus sedimentos; las
profundidades medidas incluyen los sedimentos, muchas decenas de millones de
años de ellos en un fondo así, y los montes submarinos y las mesetas
volcánicas, y todo ello hace que el fondo medido sea menos hondo.

\textbf{16.} $\alpha\,\Delta T = 3.2 \times 10^{-5} \times 700 = 0.0224$: la
parte mantélica es un \qty{2.24}{\percent} más densa,
$3380 \times 1.0224 \approx \qty{3456}{kg/m^3}$.

\textbf{17.} La parte mantélica tiene $100 - 7.1 = \qty{92.9}{km}$ de espesor.
Densidad media: $(7.1 \times 2900 + 92.9 \times 3456)/100 \approx (20590 + 321060)/100
\approx \qty{3416}{kg/m^3}$.

\textbf{18.} $3416 > 3380$: la \omterm{def:g11:lithosphere-asthenosphere:lithosphere}{litosfera} antigua en su conjunto es unos
\qty{36}{kg/m^3}, alrededor de un \qty{1}{\percent}, más densa que la
\omterm{def:g11:lithosphere-asthenosphere:lithosphere}{astenosfera} que tiene debajo. La corteza ligera ya no basta para compensar su
parte mantélica, fría y densa.

\textbf{19.} Como es más densa que la \omterm{def:g11:lithosphere-asthenosphere:lithosphere}{astenosfera}, puede volver a hundirse en
el manto en una fosa: la subducción, el tema del capítulo siguiente.

\textbf{20.} A los 64 millones de años, el trozo de fondo está a unos
\qty{5300}{m}: a unos \textbf{5.3~km} de profundidad, \qty{2.8}{km} más hondo
que donde nació.
\end{solution}
